\documentclass[../../../include/open-logic-section]{subfiles}

\begin{document}

\olfileid[id]{sth}{choice}{vitali}
\olsection{Lampiran: Paradoks Vitali}

Untuk benar-benar memahami apakah konstruksi Banach--Tarski dapat diterima,
kita perlu memeriksa \emph{buktinya}. Sayangnya, hal itu memerlukan jauh
lebih banyak aljabar daripada yang dapat kita sajikan di sini. Namun, kita
dapat memberikan beberapa catatan singkat yang mungkin memperjelas bukti
Banach--Tarski,\footnote{Untuk pembahasan yang jauh lebih lengkap, lihat
\cite{Weston2003} atau \cite{Wagon2016}.} dengan berfokus pada hasil berikut:

\begin{thm}[Paradoks Vitali (dalam $\ZFC$)]\ollabel{vitaliparadox}
Sebarang lingkaran dapat diuraikan menjadi sejumlah terhitung kepingan yang
dapat disusun kembali (melalui rotasi dan translasi) untuk membentuk dua
salinan lingkaran tersebut.
\end{thm}

Paradoks Vitali jauh lebih mudah dibuktikan daripada Paradoks Banach--Tarski.
Kita menyebutnya ``Paradoks Vitali'' karena ia mengikuti konstruksi suatu
himpunan tak terukur oleh Vitali pada \citeyear{Vitali1905}. Namun, aspek
teori-himpunan dari bukti Paradoks Vitali dan Paradoks Banach--Tarski sangat
serupa. Perbedaan mendasar antara kedua hasil itu hanyalah bahwa
Banach--Tarski mempertimbangkan suatu penguraian \emph{hingga}, sedangkan
Paradoks Vitali mempertimbangkan suatu penguraian \emph{tak hingga
terhitung}. Sebagaimana dikatakan \citet{Weston2003}, Paradoks Vitali
``tentu tidak semencolok paradoks Banach--Tarski, tetapi ia memang
menunjukkan bahwa paradoks geometris dapat terjadi bahkan dalam situasi
yang sederhana.''

Paradoks Vitali menyangkut suatu bangun dua dimensi, yakni lingkaran. Jadi,
kita akan bekerja pada bidang $\Real^2$. Misalkan $\rotationsgroup$ himpunan
rotasi titik (searah jarum jam) mengelilingi titik asal dengan besar sudut
radian \emph{rasional} dalam interval $[0,2\pi)$. Berikut beberapa fakta
aljabar mengenai $\rotationsgroup$ (jika Anda tidak memahami pernyataan
hasilnya, bukti akan memperjelas maknanya):

\begin{lem}\ollabel{rotationsgroupabelian}
$\rotationsgroup$ membentuk suatu grup abelian di bawah komposisi fungsi.
\end{lem}

\begin{proof}
Dengan menulis $0_{\rotationsgroup}$ untuk rotasi sebesar $0$ radian, rotasi
ini merupakan unsur identitas bagi $\rotationsgroup$, sebab
$\comp{0_{\rotationsgroup}}{\rho} = \comp{\rho}{0_{\rotationsgroup}} =
\rho$ bagi sebarang $\rho \in \rotationsgroup$.

Setiap anggota mempunyai invers. Jika $\rho \in \rotationsgroup$ berotasi
sebesar $r$ radian, maka $\rho^{-1} \in \rotationsgroup$ berotasi sebesar
$(-r)\bmod 2\pi$ radian, sehingga
$\rho \circ \rho^{-1} = 0_\rotationsgroup$.

Komposisi bersifat asosiatif:
$\comp{\rho}{(\comp{\sigma}{\tau})} =
\comp{(\comp{\rho}{\sigma})}{\tau}$ bagi sebarang
$\rho,\sigma,\tau\in\rotationsgroup$.

Komposisi bersifat komutatif:
$\comp{\rho}{\sigma}=\comp{\sigma}{\rho}$ bagi sebarang
$\rho,\sigma\in\rotationsgroup$.
\end{proof}

Bahkan, kita dapat membelah grup $\rotationsgroup$ menjadi dua, lalu
menggunakan salah satu belahan untuk memulihkan seluruh grup:

\begin{lem}\ollabel{disjointgroup}
Terdapat suatu partisi $\rotationsgroup$ menjadi dua himpunan saling lepas,
$\rotationsgroup_{1}$ dan $\rotationsgroup_{2}$, yang masing-masing
merupakan basis bagi $\rotationsgroup$.
\end{lem}

\begin{proof}
Misalkan $\rotationsgroup_{1}$ terdiri atas rotasi dengan besar sudut radian
rasional dalam $[0,\pi)$; misalkan
$\rotationsgroup_{2}=\rotationsgroup\setminus\rotationsgroup_1$. Berdasarkan
aljabar elementer,
$\Setabs{\comp{\rho}{\rho}}{\rho\in\rotationsgroup_1}=\rotationsgroup$.
Hasil serupa dapat diperoleh bagi $\rotationsgroup_2$.
\end{proof}

Kita akan menggunakan fakta tentang grup ini untuk menetapkan
\olref{vitaliparadox}. Misalkan $\onesphere$ lingkaran satuan, yakni himpunan
titik yang berjarak tepat $1$ satuan dari titik asal bidang, yaitu
$\Setabs{\tuple{r,s}\in\Real^2}{\sqrt{r^2+s^2}=1}$. Kita akan membagi
$\onesphere$ menjadi beberapa bagian dengan mempertimbangkan relasi berikut
pada $\onesphere$:
\[
	r \sim s \emph{ jika dan hanya jika }
	(\exists \rho \in \rotationsgroup)\rho(r) = s.
\]
Artinya, titik-titik pada $\onesphere$ dihubungkan oleh relasi ini jika dan
hanya jika kita dapat berpindah dari satu titik ke titik lain melalui rotasi
bernilai rasional mengelilingi titik asal. Tidak mengejutkan bahwa:

\begin{lem}
$\sim$ merupakan suatu relasi ekuivalensi.
\end{lem}

\begin{proof}
Trivial, dengan menggunakan \olref{rotationsgroupabelian}.
\end{proof}

Sekarang kita menggunakan Pilihan untuk memperoleh suatu himpunan $C$ yang
memuat tepat satu anggota dari setiap kelas ekuivalensi $\onesphere$ di bawah
$\sim$. Artinya, kita mempertimbangkan suatu fungsi pilihan $f$ pada
himpunan kelas-kelas ekuivalensi,\footnote{Karena $\rotationsgroup$
!!{enumerable}, setiap !!{element} dari $E$ juga !!{enumerable}. Karena
$\onesphere$ !!{nonenumerable}, berdasarkan \olref{vitalicover} dan
\olref[card-arithmetic][simp]{kappaunionkappasize}, $E$ juga
!!{nonenumerable}. Jadi, ini merupakan penggunaan Pilihan \emph{tak}
terhitung.}
\[
	E = \Setabs{\equivrep{r}{\sim}}{r \in \onesphere},
\]
dan menetapkan $C=\ran{f}$. Bagi setiap rotasi
$\rho\in\rotationsgroup$, himpunan $\funimage{\rho}{C}$ terdiri atas
titik-titik yang diperoleh dengan menerapkan rotasi $\rho$ kepada setiap
titik dalam~$C$. Dua hasil berikut menunjukkan bahwa himpunan-himpunan ini
menutupi lingkaran secara lengkap dan tanpa tumpang tindih:

\begin{lem}\ollabel{vitalicover}
$\onesphere = \bigcup_{\rho \in \rotationsgroup} \funimage{\rho}{C}$.
\end{lem}

\begin{proof}
Tetapkan $s\in\onesphere$; terdapat suatu $r\in C$ sedemikian sehingga
$r\in\equivrep{s}{\sim}$, yakni $r\sim s$, yakni
$\rho(r)=s$ bagi suatu $\rho\in\rotationsgroup$.
\end{proof}

\begin{lem}\ollabel{vitalinooverlap}
Jika $\rho_1 \neq \rho_2$, maka
$\funimage{\rho_1}{C} \cap \funimage{\rho_2}{C} = \emptyset$.
\end{lem}

\begin{proof}
Andaikan
$s\in\funimage{\rho_{1}}{C}\cap\funimage{\rho_{2}}{C}$. Jadi,
$s=\rho_{1}(r_{1})=\rho_{2}(r_{2})$ bagi suatu $r_{1},r_{2}\in C$.
Karena itu, $\rho^{-1}_2(\rho_1(r_1))=r_2$, dan
$\comp{\rho_1}{\rho^{-1}_2}\in\rotationsgroup$, sehingga
$r_1\sim r_2$. Jadi, $r_1=r_2$, sebab $C$ memilih tepat satu anggota dari
setiap kelas ekuivalensi di bawah $\sim$. Maka
$s=\rho_1(r_1)=\rho_2(r_1)$, sehingga $\rho_1=\rho_2$.
\end{proof}

Sekarang kita menerapkan fakta-fakta aljabar sebelumnya pada lingkaran kita:

\begin{lem}\ollabel{pseudobanachtarski}
Terdapat suatu partisi $\onesphere$ menjadi dua himpunan saling lepas,
$D_{1}$ dan $D_{2}$, sedemikian sehingga $D_{1}$ dapat dipartisi menjadi
sejumlah terhitung himpunan yang dapat dirotasi untuk membentuk suatu salinan
$\onesphere$ (dan demikian pula bagi $D_{2}$).
\end{lem}

\begin{proof}
Dengan menggunakan $\rotationsgroup_{1}$ dan $\rotationsgroup_{2}$ dari
\olref{disjointgroup}, tetapkan:
\begin{align*}
	D_{1} &= \bigcup_{\rho \in \rotationsgroup_1} \funimage{\rho}{C} &
	D_{2} &= \bigcup_{\rho \in \rotationsgroup_2} \funimage{\rho}{C}
\end{align*}
Ini merupakan suatu partisi $\onesphere$ berdasarkan
\olref{vitalicover}, dan $D_1$ serta $D_2$ saling lepas berdasarkan
\olref{vitalinooverlap}. Berdasarkan konstruksi, $D_1$ dapat dipartisi
menjadi sejumlah terhitung himpunan $\funimage{\rho}{C}$, satu bagi setiap
$\rho\in\rotationsgroup_1$. Himpunan-himpunan ini dapat dirotasi untuk
membentuk suatu salinan $\onesphere$, sebab
$\onesphere=\bigcup_{\rho\in\rotationsgroup}\funimage{\rho}{C}
=\bigcup_{\rho\in\rotationsgroup_1}
\funimage{(\comp{\rho}{\rho})}{C}$ berdasarkan \olref{disjointgroup} dan
\olref{vitalicover}. Penalaran yang sama berlaku bagi~$D_2$.
\end{proof}
\noindent
Hal ini langsung mengakibatkan Paradoks Vitali. Kita dapat menghasilkan
\emph{dua} salinan $\onesphere$ dari $\onesphere$ hanya dengan membaginya
menjadi sejumlah terhitung kepingan (berbagai
$\funimage{\rho}{C}$) lalu menggerakkannya secara kaku (cukup rotasikan
setiap kepingan $D_1$, serta translasikan lalu rotasikan setiap kepingan
$D_2$).

Mari kita rangkum strategi buktinya. Kita mulai dengan beberapa fakta aljabar
tentang grup rotasi pada bidang. Kita menggunakan grup ini untuk mempartisi
$\onesphere$ menjadi kelas-kelas ekuivalensi. Kemudian kita memperoleh suatu
``paradoks'' dengan menggunakan Pilihan untuk memilih anggota dari setiap
kelas.

Kita menggunakan strategi yang persis sama untuk membuktikan Banach--Tarski.
Perbedaan utamanya ialah bahwa fakta-fakta aljabar yang digunakan untuk
membuktikan Banach--Tarski jauh lebih rumit daripada yang digunakan untuk
membuktikan Paradoks Vitali. Namun, fakta-fakta aljabar tersebut tidak
berhubungan dengan Pilihan. Kita akan merangkumnya dengan cepat.

Untuk membuktikan Banach--Tarski, kita mulai dengan menetapkan suatu analog
dari \olref{disjointgroup}: sebarang \emph{grup bebas} dapat dibagi menjadi
empat kepingan yang secara intuitif dapat kita ``gerakkan'' untuk
memulihkan dua salinan seluruh grup.\footnote{Fakta bahwa kita dapat
menggunakan \emph{empat} kepingan berasal dari \cite{Robinson1947}. Untuk
suatu bukti baru-baru ini, lihat \citet[Teorema 5.2]{Wagon2016}. Kita
mengikuti \citet[hlm.~3]{Weston2003} dengan menyebutnya sebagai
``menggerakkan'' kepingan-kepingan grup.} Kemudian kita menunjukkan bahwa
kita dapat menggunakan dua rotasi tertentu mengelilingi titik asal
$\Real^3$ untuk menghasilkan suatu grup bebas rotasi, $F$.\footnote{Lihat
\citet[Teorema 2.1]{Wagon2016}.} (Belum ada Pilihan.) Sekarang kita
menganggap titik-titik pada permukaan bola ``serupa'' jika dan hanya jika
salah satunya dapat diperoleh dari yang lain melalui suatu rotasi dalam~$F$.
Kemudian kita \emph{menggunakan Pilihan} untuk memilih tepat satu titik dari
setiap kelas ekuivalensi titik-titik ``serupa''. Dengan menerapkan
pembagian $F$ pada permukaan bola seperti dalam
\olref{pseudobanachtarski}, kita membagi permukaan itu menjadi empat
kepingan yang dapat kita ``gerakkan'' untuk memperoleh dua salinan
permukaan bola. Hal ini menetapkan \citep{Hausdorff1914}:

\begin{thm}[Paradoks Hausdorff (dalam $\ZFC$)]
Permukaan sebarang bola dapat diuraikan menjadi sejumlah hingga kepingan yang
dapat disusun kembali (melalui rotasi dan translasi) untuk membentuk dua
salinan saling lepas dari bola tersebut.
\end{thm}

Beberapa trik aljabar tambahan diperlukan untuk memperoleh Teorema
Banach--Tarski sepenuhnya (yang menyangkut bukan hanya permukaan bola,
melainkan juga bagian dalamnya). Terus terang, hal ini hanya merupakan
hiasan di atas kue aljabar. Karena itu, Weston menulis:
\begin{quote}
  [\ldots] hasil tentang grup bebas merupakan \emph{langkah kunci} dalam
  bukti paradoks Banach--Tarski. Dari sudut pandang ini, paradoks
  Banach--Tarski bukanlah pernyataan mengenai $\Real^3$ sebagaimana ia
  merupakan pernyataan mengenai kompleksitas grup [translasi dan rotasi
  dalam $\Real^3$]. \cite[hlm.~16]{Weston2003}
\end{quote}
Artinya, apakah kita dapat memberikan suatu penguraian \emph{hingga} (seperti
dalam Banach--Tarski) atau penguraian \emph{tak hingga terhitung} (seperti
dalam Paradoks Vitali) ditentukan oleh fakta-fakta teori grup tertentu
tentang bekerja dalam dua atau tiga dimensi.

Memang, pengamatan terakhir ini sedikit merusak lelucon pada akhir
\olref[banach]{sec}. Karena bersifat dua dimensi, ``Banach--Tarski'' harus
dibagi menjadi sejumlah \emph{tak hingga} terhitung kepingan jika orang ingin
menyusun ulang kepingan-kepingan itu untuk membentuk
``Banach--Tarski Banach--Tarski''. Untuk memperbaiki leluconnya, orang
harus menulis dalam tiga dimensi. Kita serahkan hal ini sebagai latihan
kepada pembaca.

Satu komentar terakhir. Dalam \olref[banach]{sec}, kita menyebutkan bahwa
``kepingan-kepingan'' bola yang diperoleh tidak dapat \emph{terukur},
melainkan harus merupakan ``hamburan-hamburan tak hingga'' yang tidak
dapat digambarkan. Hal yang sama berlaku pada penggunaan Pilihan untuk
memperoleh \olref{pseudobanachtarski}. Pokok ini patut dijelaskan.

Sekali lagi, kita harus menguraikan sedikit latar belakang (tetapi ini
\emph{hanya} sebuah uraian ringkas; Anda mungkin ingin membaca entri buku
teks tentang \emph{ukuran}). Mendefinisikan ukuran bagi suatu himpunan $X$
berarti menetapkan suatu nilai $\mu(E)\in\Real$ kepada setiap $E$ dalam
suatu ``aljabar-$\sigma$'' pada~$X$. Perinciannya tidak penting di sini,
kecuali bahwa fungsi $\mu$ harus memenuhi prinsip aditivitas terhitung:
ukuran suatu gabungan terhitung dari himpunan-himpunan saling lepas merupakan
jumlah ukuran masing-masing, yakni
$\mu(\bigcup_{n < \omega} X_n) = \sum_{n < \omega}\mu(X_n)$ jika
himpunan-himpunan $X_n$ saling lepas. Mengatakan bahwa suatu himpunan
``tak terukur'' berarti tidak ada ukuran yang dapat ditetapkan kepadanya
secara semestinya. Sekarang, dengan menggunakan $\rotationsgroup$ di atas:

\begin{cor}[Vitali]
Misalkan $\mu$ suatu ukuran sedemikian sehingga $\mu(\onesphere)=1$ dan
$\mu(X)=\mu(Y)$ jika $X$ dan $Y$ kongruen. Maka
$\funimage{\rho}{C}$ tak terukur bagi semua
$\rho\in\rotationsgroup$.
\end{cor}

\begin{proof}
Untuk reductio, andaikan sebaliknya. Misalkan
$\mu(\funimage{\sigma}{C})=r$ bagi suatu
$\sigma\in\rotationsgroup$ dan suatu $r\in\Real$. Bagi sebarang
$\rho\in\rotationsgroup$, $\funimage{\rho}{C}$ dan
$\funimage{\sigma}{C}$ kongruen; karena itu,
$\mu(\funimage{\rho}{C})=r$ bagi sebarang
$\rho\in\rotationsgroup$. Berdasarkan \olref{vitalicover} dan
\olref{vitalinooverlap},
$\onesphere=\bigcup_{\rho\in\rotationsgroup}\funimage{\rho}{C}$ merupakan
suatu gabungan terhitung dari himpunan-himpunan saling lepas berpasangan.
Jadi, aditivitas terhitung mengharuskan bahwa $\mu(\onesphere)=1$ merupakan
jumlah ukuran setiap $\funimage{\rho}{C}$, yakni:
\[
	1 = \mu(\onesphere) =
	\sum_{\rho \in \rotationsgroup}\mu(\funimage{\rho}{C}) =
	\sum_{\rho \in \rotationsgroup}r
\]
Namun, jika $r=0$, maka
$\sum_{\rho\in\rotationsgroup}r=0$; dan jika $r>0$, maka
$\sum_{\rho\in\rotationsgroup}r=\infty$.
\end{proof}

\end{document}
